\section{Quiz}

% Member 4 lists the Kahoot/Quizizz multiple-choice questions here...
\label{sec:quiz}

This section evaluates the comprehension of various shortest path algorithms, their complexities, and appropriate use cases.

\begin{enumerate}[label=\arabic*., itemsep=1.5em]

    % Question 1
    \item \textbf{Question:} In the \textbf{A* search algorithm's} evaluation function $f(n) = g(n) + h(n)$, what exactly does $h(n)$ represent?
    \begin{enumerate}[label=\Alph*.]
        \item The exact accumulated cost from the start node to node $n$.
        \item The total number of edges traversed so far.
        \item The estimated cost from node $n$ to the destination.
        \item The time complexity of evaluating the current node.
    \end{enumerate}
    \textbf{Answer:} C \\
    \textbf{Explanation:} In the A* algorithm, $g(n)$ represents the exact path cost from the starting node to the current node $n$. Meanwhile, $h(n)$ is the heuristic function that estimates the optimal cost from node $n$ to the goal. This combination enables A* to prioritize nodes that are likely to lead to the most efficient path.

    % Question 2
    \item \textbf{Question:} Why does \textbf{Dijkstra's algorithm} fail to find the correct shortest path in graphs with negative edge weights?
    \begin{enumerate}[label=\Alph*.]
        \item The priority queue data structure cannot process negative integer values.
        \item It greedily marks nodes as "visited" under the assumption that path costs can never decrease.
        \item Negative edges immediately crash the algorithm due to division by zero errors.
        \item It relies on a heuristic function that strictly requires positive inputs.
    \end{enumerate}
    \textbf{Answer:} B \\
    \textbf{Explanation:} Dijkstra's algorithm fundamentally relies on the greedy property that adding an edge to a path can only increase or maintain its total length. Once a node is extracted from the priority queue and marked as visited, its shortest distance is finalized. Negative weights violate this assumption, as a path could theoretically become shorter later in the traversal.

    % Question 3
    \item \textbf{Question:} What are the standard time and space complexities of the \textbf{Floyd-Warshall algorithm}, respectively (where $V$ is the number of vertices and $E$ is the number of edges)?
    \begin{enumerate}[label=\Alph*.]
        \item Time: $\mathcal{O}(V^2)$, Space: $\mathcal{O}(V^3)$
        \item Time: $\mathcal{O}(V^3)$, Space: $\mathcal{O}(V^2)$
        \item Time: $\mathcal{O}(E \log V)$, Space: $\mathcal{O}(V)$
        \item Time: $\mathcal{O}(V^2 \log E)$, Space: $\mathcal{O}(V^2)$
    \end{enumerate}
    \textbf{Answer:} B \\
    \textbf{Explanation:} The algorithm uses three nested loops to iterate through all possible intermediate vertices for every pair of source and destination vertices, leading to a time complexity of $\mathcal{O}(V^3)$. Additionally, it requires a 2D distance matrix of size $V \times V$ to store the shortest paths, resulting in a space complexity of $\mathcal{O}(V^2)$.

    % Question 4
    \item \textbf{Question:} A graph has negative edges. To use Dijkstra's algorithm, a student adds a constant $C$ to \textbf{ALL} edges to make them non-negative ($\ge 0$), runs Dijkstra's, and subtracts the accumulated constant later. Is this method correct or incorrect, and why?
    \begin{enumerate}[label=\Alph*.]
        \item Excellent, absolutely correct. This is a classic graph transformation technique that leverages Dijkstra's speed without distorting the final result.
        \item Incorrect. Altering the weights by adding a constant will introduce hidden negative cycles into the graph, causing Dijkstra's algorithm to fall into an infinite loop.
        \item Acceptable, but the graph must be strictly undirected. If it is a directed graph, adding a constant will reverse the traversal direction.
        \item A fundamental flaw. Adding a constant $C$ to every edge heavily penalizes paths made up of a larger number of edges, completely invalidating the calculation.
    \end{enumerate}
    \textbf{Answer:} D \\
    \textbf{Explanation:} Suppose Path 1 consists of 1 edge (weight 10), and Path 2 consists of 3 edges (total weight 6). Originally, Path 2 is the shortest. If we add $C = 5$ to all edges: Path 1 becomes 15, and Path 2 becomes $6 + (3 \times 5) = 21$. Suddenly, Path 1 appears shorter. The constant addition distorts the graph's relative distances.

    % Question 5
    \item \textbf{Question:} You manage a regional delivery network. A customer requests a pickup. You need to find the most efficient route from \textit{any} of your multiple local branches to this single customer's location to dispatch a driver. \textbf{Which category of shortest-path problem are you solving?}
    \begin{enumerate}[label=\Alph*.]
        \item Single-Source Shortest Path (SSSP)
        \item Single-Pair Shortest Path (SPSP)
        \item All-Pairs Shortest Path (APSP)
        \item Multi-Source Shortest Path (MSSP)
    \end{enumerate}
    \textbf{Answer:} D \\
    \textbf{Explanation:} The scenario requires finding the optimal path from multiple potential starting points (the local branches) to a single destination. This is a classic Multi-Source problem. It can be elegantly solved by introducing a virtual "Super Source" connected to all local branches with zero-weight edges, then running an SSSP algorithm from that Super Source.

    % Question 6
    \item \textbf{Question:} Assume the current shortest path from Node A to Node B uses 5 edges, while an alternative path uses only 2 edges (but has a slightly higher total weight). If a new toll tax of $+\$10$ is added to \textit{every single edge} in the graph, \textbf{what might happen?}
    \begin{enumerate}[label=\Alph*.]
        \item The shortest path is guaranteed to remain exactly the same.
        \item The algorithm will crash because the graph's topology was altered.
        \item The alternative path (with 2 edges) might become the new shortest path.
        \item The total cost of all paths will increase by exactly $\$10$.
    \end{enumerate}
    \textbf{Answer:} C \\
    \textbf{Explanation:} Adding a uniform constant weight to all edges disproportionately impacts paths containing a larger number of edges. The cost of the original path increases by $5 \times \$10 = \$50$, whereas the alternative path only increases by $2 \times \$10 = \$20$. If the original weight difference between the two paths was less than $\$30$, the 2-edge path will now be strictly shorter.

\end{enumerate}